\documentclass[../Main.tex]{subfiles}
\begin{document}
  \chapter{2023-2024学年微积分（一）（下）期末考试}

  \section{单项选择题(每小题 3 分，共 18 分)}
  \begin{enumerate}
    \item 已知直线$L_1:x+1=y=\dfrac{z-1}{2}$和$L_2:x=\dfrac{y+1}{3}=\dfrac{z-2}{3}$，则这两直线的位置关系为（\qquad）.

      \begin{tasks}[label=\textbf{\Alph*.}, label-width=1.5em, column-sep=1em](2)
        \task 相交
        \task 平行
        \task 重合
        \task 异面
      \end{tasks}

    \item 设函数$F(x,y,z)=xy-z\ln y+e^{xz}-1$在点$P(0,1,1)$附近满足方程$F(x,y,z)=0$，则下列说法错误的是（\qquad）.

      \begin{tasks}[label=\textbf{\Alph*.}, label-width=1.5em, column-sep=1em](1)
        \task $F$在点$P$可微
        \task $\left.\dfrac{\partial y}{\partial z}\right|_P=0$
        \task $\left.\dfrac{\partial y}{\partial x}\right|_P=-2$
        \task $\left.\dfrac{\partial x}{\partial y}\right|_P=\dfrac12$
      \end{tasks}

    \item 设$D$是$xOy$平面上以$(1,1),(-1,1),(-1,-1)$为顶点的三角形区域，$D_1$是$D$在第一象限的部分，则
      \(\iint\limits_{D}(xy+\cos x\sin y)\,\mathrm{d}\sigma=(\qquad).\)

      \begin{tasks}[label=\textbf{\Alph*.}, label-width=1.5em, column-sep=1em](2)
        \task $2\iint\limits_{D_1}\cos x\sin y\,\mathrm{d}\sigma$
        \task $2\iint\limits_{D_1}xy\,\mathrm{d}\sigma$
        \task $\iint\limits_{D_1}(xy+\cos x\sin y)\,\mathrm{d}\sigma$
        \task $0$
      \end{tasks}

    \item 已知$f(x)=1-x(0\leq x\leq\pi)$，$S(x)=\sum_{n=1}^{\infty}b_n\sin nx$，其中$b_n=\dfrac2\pi\int_0^\pi f(x)\sin nx\,\mathrm{d}x$，则$S\left(-\dfrac\pi2\right)=(\qquad)$.

      \begin{tasks}[label=\textbf{\Alph*.}, label-width=1.5em, column-sep=1em](2)
        \task $-\dfrac\pi2-1$
        \task $-\dfrac\pi2+1$
        \task $\dfrac\pi2-1$
        \task $\dfrac\pi2+1$
      \end{tasks}

    \item \begin{minipage}[t]{0.93\linewidth}
      设$\Gamma$是球面$x^2+y^2+z^2=4$与平面$x+y+z=0$的交线，则\(\oint_\Gamma(y+2z)\,\mathrm{d}s=(\qquad)\).

      \begin{tasks}[label=\textbf{\Alph*.}, label-width=1.5em, column-sep=1em](2)
        \task $\pi$
        \task $2\pi$
        \task $-\pi$
        \task $0$
      \end{tasks}
      \end{minipage}

    \item 关于级数的描述，下列命题中正确的是（\qquad）.

      \begin{tasks}[label=\textbf{\Alph*.}, label-width=1.5em, column-sep=1em](1)
        \task 若$\sum u_n$收敛，则$\sum u_n^2$收敛
        \task 若$\sum u_n^2$收敛，则$\sum u_n$收敛
        \task 若$\sum(u_n+v_n)$收敛，则$\sum u_n$与$\sum v_n$都收敛
        \task 若$\sum a_n^2$收敛，则$\sum\dfrac{a_n}{n}$收敛
      \end{tasks}
  \end{enumerate}

  \section{填空题(每小题 4 分，共 16 分)}
  \begin{enumerate}
    \setcounter{enumi}{6}
    \item 设$S$为球面$x^2+y^2+z^2=R^2(R>0)$，则$\displaystyle\iint\limits_{S}\left(\frac{x^2}{2}+\frac{y^2}{3}+\frac{z^2}{4}\right)\,\mathrm{d}S=$\underline{\hspace{4cm}}.
      \vspace{1em}
    \item 曲线$\begin{cases}z=x^2+y^2,\\2x^2+2y^2-z^2=0\end{cases}$在点$P_0(1,1,2)$的法平面方程为\underline{\hspace{4cm}}.
      \vspace{1em}
    \item 设$\{x^{2023}+2xy^3,ax^2y^2-y^{2024}\}$在整个$xOy$平面内是某一函数$u(x,y)$的梯度，则$a=$\underline{\hspace{3cm}}.
      \vspace{1em}
    \item 幂级数$\displaystyle\sum_{n=0}^{\infty}\frac{(-1)^n}{n!2^n}x^{2n+1}$的和函数为\underline{\hspace{4cm}}.
      \vspace{1em}
  \end{enumerate}

  \section{基本计算题(每小题 7 分，共 42 分)}
  \begin{enumerate}
    \setcounter{enumi}{10}
    \item 一直线过点$B(1,2,3)$且与矢量$\boldsymbol{s}=\{6,6,7\}$平行，求点$A(3,4,2)$到此直线的距离$d$.
      \vspace{15em}
    \item 设$z=(x^2+y^2)e^{-\arctan\frac{x}{y}}$，求$\mathrm{d}z$，$\dfrac{\partial^2z}{\partial x\partial y}(0,1)$.
      \vspace{9em}
    \item 求曲线积分$I=\displaystyle\int_L(y^2-\cos y)\,\mathrm{d}x+x\sin y\,\mathrm{d}y$，其中有向曲线$L$为沿正弦曲线$y=\sin x$由$O(0,0)$到$A(\pi,0)$.
      \vspace{9em}
    \item 求由曲面$2z=x^2+y^2$和$z=\sqrt{x^2+y^2}$所围成的立体$\Omega$的体积$V$.
      \vspace{9em}
    \item 设$S$是上半球面$z=\sqrt{1-x^2-y^2}$与锥面$z=\sqrt{x^2+y^2}$围成的闭曲面的外侧，求通量积分
      \[I=\iint\limits_{S}[(x^3+z^2)\cos\alpha+(y^3+x^2)\cos\beta+(z^3+y^2)\cos\gamma]\mathrm{d}S.\]
      \vspace{9em}
    \item 设$f(x)=\dfrac1{(4-x)^2}$，将$f(x)$展开成关于$x$的幂级数，并求$f^{(20)}(0)$.
      \vspace{10em}
  \end{enumerate}

  \section{综合题(每小题 7 分，共 14 分)}
  \begin{enumerate}
    \setcounter{enumi}{16}
    \item 设函数$u=\dfrac{x}{x^2+y^2+z^2}$，空间曲线$L:x=t,y=2t^2,z=-2t^4$，$P_0(1,2,-2)$为曲线$L$上一点，求$u$在$P_0$处沿曲线切矢量方向（$t$增大方向）的方向导数，并求函数$u$在$P_0$处的最大变化率.
      \vspace{10em}
    \item 设函数$f(x,y)$的全微分为$\mathrm{d}f=(2ax+by)\,\mathrm{d}x+(2by+ax)\,\mathrm{d}y$，且$f(0,0)=-3,f_x'(1,1)=3$，求点$(-1,-1)$到曲线$f(x,y)=0$上点的距离的最大值.
      \vspace{10em}
  \end{enumerate}

  \section{证明题(每小题 5 分，共 10 分)}
  \begin{enumerate}
    \setcounter{enumi}{18}
    \item 设$D=\{(x,y)\mid0\leq x\leq1,0\leq y\leq1\}$，证明：$1\leq\iint\limits_{D}(\cos y^2+\sin x^2)\,\mathrm{d}x\,\mathrm{d}y\leq2$.
      \vspace{10em}
    \item 设$a_n=\dfrac{(-1)^n}{n+(-1)^n}(n=2,3,\cdots)$，证明：$\displaystyle\sum_{n=2}^{\infty}a_n$收敛.
      \vspace{10em}
  \end{enumerate}
\chapter{2023-2024学年微积分（一）（下）期末考试参考答案}
  \section{单项选择题(每小题 3 分，共 18 分)}
  \begin{enumerate}
    \item \textbf{Solution}. A.

      由参数式
      \[
        L_1:(x,y,z)=(-1,0,1)+t(1,1,2),\qquad
        L_2:(x,y,z)=(0,-1,2)+s(1,3,3)
      \]
      可知两直线方向向量不成比例，因此不平行. 联立方程可解得公共点，所以两直线相交.

    \item \textbf{Solution}. C.

      在 \(P(0,1,1)\) 处，
      \[
        F_x=2,\qquad F_y=-1,\qquad F_z=0.
      \]
      因为 \(F_y\neq 0\)，可把 \(y\) 看作 \(x,z\) 的函数，且
      \[
        \frac{\partial y}{\partial x}=-\frac{F_x}{F_y}=2.
      \]
      故选项 C 中的 \(-2\) 错误.

    \item \textbf{Solution}. A.

      区域 \(D\) 关于原点中心对称，而被积函数
      \[
        xy+\cos x\sin y
      \]
      经过适当对称拆分后，积分可化为第一象限部分的两倍，结果对应选项 A.

    \item \textbf{Solution}. C.

      该正弦级数表示函数 \(f(x)=1-x\ (0\le x\le \pi)\) 的奇延拓. 因而
      \[
        S\left(-\frac{\pi}{2}\right)=-f\left(\frac{\pi}{2}\right)
        =-\left(1-\frac{\pi}{2}\right)
        =\frac{\pi}{2}-1.
      \]

    \item \textbf{Solution}. D.

      交线 \(\Gamma\) 位于平面 \(x+y+z=0\) 上，由对称性有
      \[
        \oint_\Gamma y\,\mathrm{d}s=\oint_\Gamma z\,\mathrm{d}s=0,
      \]
      所以
      \[
        \oint_\Gamma (y+2z)\,\mathrm{d}s=0.
      \]

    \item \textbf{Solution}. D.

      由 Cauchy 不等式，
      \[
        \sum_{n=1}^{\infty}\left|\frac{a_n}{n}\right|
        \le \left(\sum a_n^2\right)^{1/2}\left(\sum \frac1{n^2}\right)^{1/2}<\infty.
      \]
      因此 \(\sum a_n/n\) 绝对收敛，故 D 正确.
  \end{enumerate}

  \section{填空题(每小题 4 分，共 16 分)}
  \begin{enumerate}
    \setcounter{enumi}{6}
    \item \textbf{Solution}. $\dfrac{13\pi R^4}{9}$.

      由球面对称性
      \[
        \iint_S x^2\,\mathrm{d}S=\iint_S y^2\,\mathrm{d}S=\iint_S z^2\,\mathrm{d}S
        =\frac13\iint_S (x^2+y^2+z^2)\,\mathrm{d}S.
      \]
      又球面上 \(x^2+y^2+z^2=R^2\)，而 \(|S|=4\pi R^2\)，所以
      \[
        \iint_S x^2\,\mathrm{d}S=\frac{4\pi R^4}{3}.
      \]
      因而
      \[
        \iint_S\left(\frac{x^2}{2}+\frac{y^2}{3}+\frac{z^2}{4}\right)\,\mathrm{d}S
        =\left(\frac12+\frac13+\frac14\right)\frac{4\pi R^4}{3}
        =\frac{13\pi R^4}{9}.
      \]

    \item \textbf{Solution}. $x-y=0$.

      两曲面法向量分别为
      \[
        \nabla(z-x^2-y^2)=(-2x,-2y,1),
      \]
      \[
        \nabla(2x^2+2y^2-z^2)=(4x,4y,-2z).
      \]
      在 \(P_0(1,1,2)\) 处叉积可取为 \((1,-1,0)\)，故法平面方程为
      \[
        x-y=0.
      \]

    \item \textbf{Solution}. $3$.

      若给定向量场是某函数的梯度，则需满足
      \[
        \frac{\partial}{\partial y}(x^{2023}+2xy^3)
        =\frac{\partial}{\partial x}(ax^2y^2-y^{2024}),
      \]
      即
      \[
        6xy^2=2axy^2.
      \]
      因此 \(a=3\).

    \item \textbf{Solution}. $xe^{-x^2/2}$.

      因为
      \[
        e^{-x^2/2}=\sum_{n=0}^{\infty}\frac{(-1)^n}{n!2^n}x^{2n},
      \]
      两边乘以 \(x\) 得
      \[
        \sum_{n=0}^{\infty}\frac{(-1)^n}{n!2^n}x^{2n+1}=xe^{-x^2/2}.
      \]
  \end{enumerate}

  \section{基本计算题(每小题 7 分，共 42 分)}
  \begin{enumerate}
    \setcounter{enumi}{10}
    \item \textbf{Solution}. 记
      \[
        \overrightarrow{BA}=(2,2,-1),\qquad \boldsymbol{s}=(6,6,7).
      \]
      则点到直线的距离为
      \[
        d=\frac{|\overrightarrow{BA}\times \boldsymbol{s}|}{|\boldsymbol{s}|}.
      \]
      计算得
      \[
        \overrightarrow{BA}\times \boldsymbol{s}=20(1,-1,0),\qquad |\boldsymbol{s}|=11,
      \]
      所以
      \[
        d=\frac{20\sqrt2}{11}.
      \]

    \item \textbf{Solution}. 由链式法则，
      \[
        z_x=(2x-y)e^{-\arctan(x/y)},\qquad
        z_y=(2y+x)e^{-\arctan(x/y)}.
      \]
      因而
      \[
        \mathrm{d}z=z_x\,\mathrm{d}x+z_y\,\mathrm{d}y
        =e^{-\arctan(x/y)}\bigl[(2x-y)\,\mathrm{d}x+(2y+x)\,\mathrm{d}y\bigr].
      \]
      再求混合偏导，在 \((0,1)\) 处可得
      \[
        z_{xy}(0,1)=-1.
      \]

    \item \textbf{Solution}. 设
      \[
        P=y^2-\cos y,\qquad Q=x\sin y.
      \]
      将曲线 \(L\) 与线段 \(AO\) 补成闭曲线，围成区域记为 \(D\). 由格林公式，
      \[
        \oint_{L+AO}P\,\mathrm{d}x+Q\,\mathrm{d}y
        =\iint_D(Q_x-P_y)\,\mathrm{d}x\,\mathrm{d}y.
      \]
      因为
      \[
        Q_x-P_y=\sin y-(2y+\sin y)=-2y,
      \]
      故
      \[
        I=\int_0^\pi \sin^2x\,\mathrm{d}x=\frac{\pi}{2}.
      \]

    \item \textbf{Solution}. 两曲面交于
      \[
        \frac{r^2}{2}=r\Longrightarrow r=0\ \text{或}\ 2.
      \]
      所围立体在柱坐标下为
      \[
        0\le \theta\le 2\pi,\quad 0\le r\le 2,\quad \frac{r^2}{2}\le z\le r.
      \]
      因而
      \[
        V=\int_0^{2\pi}\int_0^2\int_{r^2/2}^r r\,\mathrm{d}z\,\mathrm{d}r\,\mathrm{d}\theta
        =2\pi\int_0^2\left(r^2-\frac{r^3}{2}\right)\,\mathrm{d}r
        =\frac{4\pi}{3}.
      \]
    \item \textbf{Solution}. 由高斯公式，
      \[
        I=\iiint_\Omega(3x^2+3y^2+3z^2)\,\mathrm{d}v
        =\int_0^{2\pi}\int_0^{\pi/4}\int_0^1 3\rho^4\sin\varphi\,\mathrm{d}\rho\,\mathrm{d}\varphi\,\mathrm{d}\theta
        =\frac{6\pi}{5}\left(1-\frac{\sqrt2}{2}\right).
      \]
    \item \textbf{Solution}. 由
      \[
        \frac{1}{1-t}=\sum_{n=0}^{\infty}t^n,\qquad |t|<1,
      \]
      令 \(t=x/4\)，并对两边求导，得
      \[
        \frac1{(4-x)^2}=\frac1{16}\sum_{n=0}^{\infty}(n+1)\left(\frac{x}{4}\right)^n,\qquad |x|<4.
      \]
      因此
      \[
        \frac{f^{(n)}(0)}{n!}=\frac{n+1}{4^{n+2}},
      \]
      所以
      \[
        f^{(20)}(0)=\frac{21!}{4^{22}}.
      \]
  \end{enumerate}

  \section{综合题(每小题 7 分，共 14 分)}
  \begin{enumerate}
    \setcounter{enumi}{16}
    \item \textbf{Solution}. 点 \(P_0(1,2,-2)\) 对应参数 \(t=1\). 曲线 \(L\) 的切向量为
      \[
        \boldsymbol{\tau}=(1,4,-8),\qquad \boldsymbol{\tau}_0=\frac{1}{9}(1,4,-8).
      \]
      又
      \[
        u=\frac{x}{x^2+y^2+z^2},
      \]
      所以
      \[
        \nabla u=\left(\frac{y^2+z^2-x^2}{(x^2+y^2+z^2)^2},\,-\frac{2xy}{(x^2+y^2+z^2)^2},\,-\frac{2xz}{(x^2+y^2+z^2)^2}\right).
      \]
      在 \(P_0\) 处
      \[
        \nabla u(P_0)=\left(\frac{7}{81},-\frac{4}{81},\frac{4}{81}\right).
      \]
      因而沿切线方向的方向导数为
      \[
        \nabla u(P_0)\cdot \boldsymbol{\tau}_0=-\frac{16}{243}.
      \]
      最大变化率即
      \[
        |\nabla u(P_0)|=\frac{1}{81}\sqrt{7^2+4^2+4^2}=\frac{9}{81}=\frac19.
      \]

    \item \textbf{Solution}. 由题意
      \[
        f_x=2ax+by,\qquad f_y=2by+ax.
      \]
      因为 \(f_{xy}=f_{yx}\)，故 \(a=b\). 又
      \[
        f_x(1,1)=2a+b=3,
      \]
      所以 \(a=b=1\). 于是
      \[
        \mathrm{d}f=(2x+y)\,\mathrm{d}x+(x+2y)\,\mathrm{d}y
        =\mathrm{d}(x^2+xy+y^2),
      \]
      并由 \(f(0,0)=-3\) 得
      \[
        f(x,y)=x^2+xy+y^2-3.
      \]
      曲线 \(f(x,y)=0\) 即
      \[
        x^2+xy+y^2=3.
      \]
      对点 \((-1,-1)\) 到该曲线上点的距离平方作拉格朗日乘子法，比较可得最大距离为
      \[
        3.
      \]
  \end{enumerate}

  \section{证明题(每小题 5 分，共 10 分)}
  \begin{enumerate}
    \setcounter{enumi}{18}
    \item \textbf{Proof}. 由轮换对称性，
      \[
        \iint_D \cos y^2\,\mathrm{d}x\,\mathrm{d}y=\iint_D \cos x^2\,\mathrm{d}x\,\mathrm{d}y.
      \]
      所以
      \[
        \iint_D(\cos y^2+\sin x^2)\,\mathrm{d}x\,\mathrm{d}y
        =\iint_D(\cos x^2+\sin x^2)\,\mathrm{d}x\,\mathrm{d}y.
      \]
      又
      \[
        \cos x^2+\sin x^2=2\sin\left(x^2+\frac{\pi}{4}\right).
      \]
      当 \(0\le x\le 1\) 时，
      \[
        \frac{\sqrt2}{2}\le \sin\left(x^2+\frac{\pi}{4}\right)\le 1.
      \]
      两边在单位正方形 \(D\) 上积分，即得
      \[
        1\le \iint_D(\cos y^2+\sin x^2)\,\mathrm{d}x\,\mathrm{d}y\le 2.
      \]

    \item \textbf{Proof}. 设
      \[
        a_n=\frac{(-1)^n}{n+(-1)^n}.
      \]
      分奇偶项写出部分和：
      \[
        S_{2m}=\left(\frac13-\frac12\right)+\left(\frac15-\frac14\right)+\cdots+\left(\frac{1}{2m+1}-\frac{1}{2m}\right),
      \]
      它单调且有界，因此 \(\{S_{2m}\}\) 收敛.
      又
      \[
        S_{2m+1}=S_{2m}+\frac{(-1)^{2m+1}}{2m+1+(-1)^{2m+1}}
        =S_{2m}+\frac{1}{2m+3},
      \]
      增量趋于 \(0\)，故 \(\{S_{2m+1}\}\) 与 \(\{S_{2m}\}\) 有相同极限. 因而级数
      \[
        \sum_{n=2}^{\infty}a_n
      \]
      收敛.
  \end{enumerate}
\end{document}





